
\documentclass[10pt,aspectratio=169]{beamer}
\usepackage[utf8]{inputenc}
\usepackage[T1]{fontenc}
\usepackage[catalan]{babel}
\usepackage{amsmath,amssymb,booktabs,array,multicol}
\usepackage{tikz,pgfplots,xcolor,listings}
\usetikzlibrary{positioning,arrows.meta,calc,shapes.geometric}
\pgfplotsset{compat=1.18}
\usepgfplotslibrary{fillbetween}

\definecolor{UABGreen}{RGB}{0,103,71}
\definecolor{UABLight}{RGB}{224,241,236}
\definecolor{DeepBlue}{RGB}{20,65,110}
\definecolor{AccentOrange}{RGB}{230,126,34}
\definecolor{SoftGray}{RGB}{245,247,250}
\definecolor{CodeBack}{RGB}{248,248,248}
\definecolor{CodeFrame}{RGB}{210,215,220}

\usetheme{Madrid}
\usecolortheme[named=UABGreen]{structure}
\setbeamercolor{title}{fg=white,bg=UABGreen}
\setbeamercolor{frametitle}{fg=white,bg=UABGreen}
\setbeamercolor{block title}{fg=white,bg=DeepBlue}
\setbeamercolor{block body}{fg=black,bg=SoftGray}
\setbeamercolor{alerted text}{fg=AccentOrange}
\setbeamertemplate{navigation symbols}{}
\setbeamertemplate{itemize item}{\color{UABGreen}$\blacktriangleright$}
\setbeamertemplate{itemize subitem}{\color{AccentOrange}$\bullet$}
\setbeamertemplate{enumerate item}{\color{UABGreen}\insertenumlabel.}
\setbeamertemplate{footline}{%
  \leavevmode%
  \hbox{%
  \begin{beamercolorbox}[wd=.72\paperwidth,ht=2.4ex,dp=1ex,leftskip=1.4em]{author in head/foot}%
    Estadística -- Grau en Enginyeria Informàtica
  \end{beamercolorbox}%
  \begin{beamercolorbox}[wd=.28\paperwidth,ht=2.4ex,dp=1ex,center]{date in head/foot}%
    \insertframenumber{} / \inserttotalframenumber
  \end{beamercolorbox}}%
  \vskip0pt%
}

\lstdefinelanguage{R}{%
  morekeywords={if,else,repeat,while,function,for,in,next,break,TRUE,FALSE,NULL,NA,NaN,Inf,
    library,mean,median,sd,var,summary,table,prop.table,xtabs,addmargins,round,
    ggplot,aes,geom_point,geom_line,geom_col,geom_histogram,geom_density,geom_smooth,
    geom_tile,facet_wrap,labs,theme_minimal,summarise,group_by,mutate,read_csv,filter,
    arrange,set.seed,rbinom,rgeom,rpois,runif,rexp,rnorm,dnorm,pnorm,qnorm,dt,pt,qt,
    dchisq,pchisq,qchisq,df,pf,qf,dbinom,pbinom,dpois,ppois,replicate,sample,rep,seq,
    tibble,data.frame,confint,t.test,var.test,chisq.test,binom.test,prop.test,lm,
    broom,glance,augment,tidy,future_map_dbl,plan,multisession,map_dbl,purrr,future,furrr},
  sensitive=true,
  morecomment=[l]\#,
  morestring=[b]",
  morestring=[b]'
}
\lstset{language=R,basicstyle=\ttfamily\scriptsize,keywordstyle=\color{DeepBlue}\bfseries,
  commentstyle=\color{gray},stringstyle=\color{AccentOrange},backgroundcolor=\color{CodeBack},
  frame=single,rulecolor=\color{CodeFrame},breaklines=true,columns=fullflexible,keepspaces=true,
  showstringspaces=false,
  literate={à}{{\`a}}1 {è}{{\`e}}1 {é}{{\'e}}1 {í}{{\'i}}1 {ï}{{\"i}}1 {ò}{{\`o}}1 {ó}{{\'o}}1 {ú}{{\'u}}1 {ü}{{\"u}}1 {ç}{{\c{c}}}1 {À}{{\`A}}1 {È}{{\`E}}1 {É}{{\'E}}1 {Í}{{\'I}}1 {Ò}{{\`O}}1 {Ó}{{\'O}}1 {Ú}{{\'U}}1}

\newcommand{\R}{\textsf{R}}
\newcommand{\RStudio}{\textsf{RStudio/Posit}}
\newcommand{\GitHub}{\textsf{GitHub}}
\newcommand{\important}[1]{\textcolor{UABGreen}{\textbf{#1}}}
\newcommand{\exemple}[1]{\textcolor{AccentOrange}{\textbf{#1}}}
\newcommand{\Prob}{\mathbb{P}}
\newcommand{\E}{\mathbb{E}}
\newcommand{\Var}{\operatorname{Var}}
\newcommand{\IC}{\operatorname{IC}}

\title[Inferència IV]{Inferència IV}
\author{Estadística -- Grau en Enginyeria Informàtica}
\date{Curs 2026--27}

\begin{document}
\begin{frame}[plain]
  \titlepage
  \vspace{-0.45cm}
  \begin{center}
  \small\textcolor{DeepBlue}{Estadística computacional amb \R, simulació i exemples d'enginyeria informàtica}
  \end{center}
\end{frame}


\begin{frame}[fragile]{Per què comparar mitjanes?}
\begin{itemize}
  \item Volem saber si dues configuracions tenen el mateix rendiment mitjà.
  \item Comparar una versió nova amb una versió anterior és una situació habitual.
  \item La pregunta estadística és si la diferència observada supera el soroll mostral.
\end{itemize}
\begin{block}{Exemples}
Latència abans/després, temps d'execució de dos algoritmes, consum de memòria amb dues configuracions.
\end{block}
\end{frame}

\begin{frame}[fragile]{Objectius}
\begin{enumerate}
  \item Distingir dades aparellades i mostres independents.
  \item Formular tests de comparació de mitjanes.
  \item Aplicar tests t adequats amb \R.
  \item Interpretar p-valors, intervals i diferències estimades.
  \item Relacionar-ho amb benchmarking i experiments A/B quantitatius.
\end{enumerate}
\end{frame}

\begin{frame}[fragile]{Dades aparellades o independents?}
\begin{columns}[T]
\column{0.48\textwidth}
\begin{block}{Aparellades}
Les dues mesures corresponen a la mateixa unitat.
\begin{itemize}
  \item mateix servidor abans/després,
  \item mateix conjunt de tasques amb dos algoritmes,
  \item mateix usuari en dues condicions.
\end{itemize}
\end{block}
\column{0.48\textwidth}
\begin{block}{Independents}
Les dues mostres provenen de grups diferents.
\begin{itemize}
  \item dos grups d'usuaris,
  \item dues màquines diferents,
  \item dues mostres de peticions no relacionades.
\end{itemize}
\end{block}
\end{columns}
\end{frame}

\begin{frame}[fragile]{Dades aparellades: reduir a un test simple}
\begin{block}{Diferències}
Si tenim parelles $(X_i,Y_i)$, definim
\[
D_i=Y_i-X_i.
\]
El test sobre dues mitjanes es converteix en un test sobre una mitjana:
\[
H_0:\mu_D=0.
\]
\end{block}
\begin{lstlisting}
abans  <- c(120,118,130,125,119,122,121,128)
despres <- c(112,116,124,120,116,119,117,121)
t.test(despres, abans, paired = TRUE,
       alternative = "less")
\end{lstlisting}
\end{frame}

\begin{frame}[fragile]{Mostres independents: test de Welch}
\begin{block}{Opció recomanada per defecte}
El test t de Welch no exigeix variàncies iguals:
\[
H_0:\mu_1=\mu_2,\qquad H_1:\mu_1\ne\mu_2.
\]
\end{block}
\begin{lstlisting}
alg_A <- c(34, 31, 33, 36, 35, 32, 37, 34)
alg_B <- c(29, 30, 31, 28, 32, 30, 29, 31)
t.test(alg_A, alg_B, alternative = "greater")
\end{lstlisting}
\begin{itemize}
  \item \texttt{greater}: la mitjana d'A és més gran que la de B.
  \item En temps d'execució, ``més gran'' pot voler dir pitjor.
\end{itemize}
\end{frame}

\begin{frame}[fragile]{Mostres independents amb variàncies iguals}
Si assumim normalitat i variàncies iguals, es pot usar la versió amb variància combinada.
\begin{lstlisting}
t.test(alg_A, alg_B,
       alternative = "greater",
       var.equal = TRUE)
\end{lstlisting}
\begin{alertblock}{Prudència}
En dades de rendiment, les variàncies sovint poden diferir entre configuracions. Welch és més robust.
\end{alertblock}
\end{frame}

\begin{frame}[fragile]{Exemple: comparació d'algoritmes}
\begin{columns}[T]
\column{0.52\textwidth}
Tenim temps d'execució de dos algoritmes sobre problemes independents.
\begin{itemize}
  \item $H_0:\mu_A=\mu_B$.
  \item $H_1:\mu_A>\mu_B$ si volem demostrar que A és més lent.
  \item També podem plantejar $H_1:\mu_A<\mu_B$ si volem demostrar que A és més ràpid.
\end{itemize}
\column{0.42\textwidth}
\begin{block}{Informe}
No direm només ``p < 0.05''. Cal indicar diferència estimada i interval de confiança.
\end{block}
\end{columns}
\end{frame}

\begin{frame}[fragile]{Visualització abans del test}
\begin{lstlisting}
library(tidyverse)
dades <- tibble(
  temps = c(alg_A, alg_B),
  algoritme = rep(c("A", "B"), each = length(alg_A))
)

ggplot(dades, aes(algoritme, temps)) +
  geom_boxplot() +
  geom_point() +
  theme_minimal()
\end{lstlisting}
\begin{block}{Regla}
Abans d'un test, sempre visualitzem les dades.
\end{block}
\end{frame}

\begin{frame}[fragile]{Criteris pràctics de decisió}
\begin{itemize}
  \item Estadística: p-valor i interval de confiança.
  \item Enginyeria: diferència rellevant per a l'usuari o el sistema.
  \item Cost: complexitat afegida, manteniment, estabilitat.
  \item Reproduïbilitat: scripts, llavors, entorn i dades versionades.
\end{itemize}
\begin{alertblock}{Pregunta final}
La diferència és estadísticament detectable i pràcticament útil?
\end{alertblock}
\end{frame}

\end{document}
